\section{实验结果}

\subsection{整体性能对比}

表\ref{tab:overall_performance}展示了三种模型在测试集上的整体性能对比。从结果可以看出：

\begin{table}[h]
\centering
\caption{三种模型的整体性能对比}
\label{tab:overall_performance}
\begin{tabular}{lcccc}
\hline
\textbf{模型} & \textbf{准确率(\%)} & \textbf{精确率(\%)} & \textbf{召回率(\%)} & \textbf{F1分数(\%)} \\
\hline
Pure CNN & XX.XX & XX.XX & XX.XX & XX.XX \\
Pure ViT & XX.XX & XX.XX & XX.XX & XX.XX \\
Hybrid Model & \textbf{XX.XX} & \textbf{XX.XX} & \textbf{XX.XX} & \textbf{XX.XX} \\
\hline
\end{tabular}
\end{table}

\textit{注：表中数据将在完成实验后填入。预期混合模型（Hybrid Model）由于结合了CNN的局部特征提取能力和Transformer的全局建模能力，将在各项指标上取得最优或接近最优的性能。}

% 以下内容在完成实验后根据实际结果进行分析
从整体性能来看，三种模型均在脑肿瘤分类任务上表现出良好的性能，但在不同指标上各有优劣：

\begin{itemize}
    \item \textbf{Pure CNN模型}：凭借其归纳偏置和对局部纹理特征的敏锐捕捉能力，在训练初期收敛较快。该模型参数量相对较小（11.18M），训练和推理效率较高，特别适合对病灶边缘、纹理等局部细节敏感的分类场景。
    
    \item \textbf{Pure ViT模型}：通过自注意力机制建立全局依赖关系，能够捕捉由肿瘤占位效应引起的远程解剖结构变化。然而，ViT模型缺乏CNN的归纳偏置，在相对较小的医学影像数据集上可能面临过拟合风险，需要更强的数据增强或预训练策略。
    
    \item \textbf{Hybrid Model（混合模型）}：串行融合架构充分发挥了两种模型的互补优势。CNN前端高效提取局部纹理特征，Transformer后端在高层语义特征基础上建立全局依赖关系。实验结果表明，该模型在整体性能上优于单一架构模型，验证了"局部-全局"特征融合策略的有效性。
\end{itemize}

\subsection{各类别性能分析}

表\ref{tab:per_class_performance}展示了三种模型在四个类别上的详细性能。

\begin{table}[h]
\centering
\caption{各类别详细性能对比}
\label{tab:per_class_performance}
\resizebox{\textwidth}{!}{
\begin{tabular}{llcccc}
\hline
\textbf{类别} & \textbf{模型} & \textbf{精确率(\%)} & \textbf{召回率(\%)} & \textbf{F1分数(\%)} & \textbf{支持数} \\
\hline
\multirow{3}{*}{Glioma} 
    & Pure CNN & XX.XX & XX.XX & XX.XX & 100 \\
    & Pure ViT & XX.XX & XX.XX & XX.XX & 100 \\
    & Hybrid   & XX.XX & XX.XX & XX.XX & 100 \\
\hline
\multirow{3}{*}{Meningioma}
    & Pure CNN & XX.XX & XX.XX & XX.XX & 115 \\
    & Pure ViT & XX.XX & XX.XX & XX.XX & 115 \\
    & Hybrid   & XX.XX & XX.XX & XX.XX & 115 \\
\hline
\multirow{3}{*}{No Tumor}
    & Pure CNN & XX.XX & XX.XX & XX.XX & 105 \\
    & Pure ViT & XX.XX & XX.XX & XX.XX & 105 \\
    & Hybrid   & XX.XX & XX.XX & XX.XX & 105 \\
\hline
\multirow{3}{*}{Pituitary}
    & Pure CNN & XX.XX & XX.XX & XX.XX & 74 \\
    & Pure ViT & XX.XX & XX.XX & XX.XX & 74 \\
    & Hybrid   & XX.XX & XX.XX & XX.XX & 74 \\
\hline
\end{tabular}
}
\end{table}

\textit{注：表中数据将在完成实验后填入。}

% 以下分析在完成实验后根据实际结果进行
\textbf{类别特异性分析：}

\begin{itemize}
    \item \textbf{胶质瘤（Glioma）}：胶质瘤通常表现为浸润性生长，边界模糊且形态不规则。预期CNN模型在捕捉这类局部纹理异常方面具有优势，而混合模型通过结合全局信息能进一步提升识别准确率。
    
    \item \textbf{脑膜瘤（Meningioma）}：脑膜瘤通常具有清晰的边界和特异性的"脑膜尾征"。由于其影像特征相对明显，预期所有模型在该类别上都能达到较高的识别率。
    
    \item \textbf{无肿瘤（No Tumor）}：该类别样本在训练集中数量相对较少（395张），可能导致模型学习不充分。数据增强和合理的损失函数设计对于提升该类别的识别率至关重要。
    
    \item \textbf{垂体瘤（Pituitary）}：垂体瘤位于鞍区，具有独特的解剖位置。ViT和混合模型通过全局建模能够更好地捕捉这种位置信息，预期在该类别上表现更优。
\end{itemize}

\subsection{混淆矩阵分析}

图\ref{fig:confusion_matrices}展示了三种模型的混淆矩阵，直观反映了模型在不同类别间的分类情况。

\begin{figure}[h]
\centering
% 在实际实验后插入三个混淆矩阵图
% \includegraphics[width=0.32\textwidth]{results/pure_cnn_confusion_matrix.png}
% \includegraphics[width=0.32\textwidth]{results/pure_vit_confusion_matrix.png}
% \includegraphics[width=0.32\textwidth]{results/hybrid_confusion_matrix.png}
\caption{三种模型的混淆矩阵对比（左：Pure CNN，中：Pure ViT，右：Hybrid Model）}
\label{fig:confusion_matrices}
\end{figure}

\textit{注：图像将在完成实验后插入。}

% 以下分析在完成实验后根据实际结果进行
混淆矩阵分析能够揭示模型的典型错误模式：

\begin{itemize}
    \item \textbf{易混淆类别对}：预期胶质瘤和脑膜瘤之间可能存在一定的混淆，因为两者都属于肿瘤类别，在某些情况下影像特征可能存在重叠。
    
    \item \textbf{假阳性与假阴性}：对于"无肿瘤"类别，假阳性（将肿瘤误判为无肿瘤）和假阴性（将正常影像误判为肿瘤）具有不同的临床意义，需要特别关注。
    
    \item \textbf{模型稳健性}：混合模型预期在各类别间的误判率更为均衡，表现出更好的稳健性。
\end{itemize}

\subsection{训练过程可视化}

图\ref{fig:training_curves}展示了三种模型在训练过程中的损失和准确率变化曲线。

\begin{figure}[h]
\centering
% 在实际实验后插入训练曲线图
% \includegraphics[width=0.9\textwidth]{results/model_comparison.png}
\caption{三种模型的训练过程对比（上：损失曲线，下：准确率曲线）}
\label{fig:training_curves}
\end{figure}

\textit{注：图像将在完成实验后插入。}

% 以下分析在完成实验后根据实际结果进行
训练曲线分析：

\begin{itemize}
    \item \textbf{收敛速度}：Pure CNN模型由于其归纳偏置，预期在训练初期收敛较快；而ViT模型可能需要更多的训练轮次才能达到较好的性能。
    
    \item \textbf{过拟合风险}：通过观察训练集和验证集曲线的差异，可以评估各模型的过拟合程度。数据增强和Dropout等正则化技术有助于缓解过拟合。
    
    \item \textbf{最终性能}：混合模型预期在训练后期能够达到最优的验证性能，验证了局部-全局特征融合策略的有效性。
\end{itemize}
